\documentclass[11pt]{article}
\usepackage[a4paper,margin=1in]{geometry}
\usepackage{fontspec}
\usepackage[boldfont=false]{xeCJK}
\setCJKmainfont{FandolSong-Regular.otf}
\usepackage{amsmath}
\usepackage{amssymb}
\usepackage{array}
\usepackage{booktabs}
\usepackage{graphicx}
\usepackage{adjustbox}
\usepackage{enumitem}
\setlist[itemize]{topsep=2pt,partopsep=0pt,parsep=0pt,itemsep=2pt,leftmargin=1.4em}
\setlist[enumerate]{topsep=2pt,partopsep=0pt,parsep=0pt,itemsep=2pt,leftmargin=1.6em}
\usepackage{parskip}
\usepackage{fvextra}
\fvset{breaklines=true,breakanywhere=true,fontsize=\small}
\setlength{\parindent}{0pt}

\begin{document}

\begin{center}
  {\LARGE\bfseries Gravitational fields}\\[0.35em]
  {\large A Level Physics}
\end{center}
\vspace{0.6em}

\section*{Gravitational fields}

\subsection*{Definition}

A \textbf{gravitational field} 重力场 is a region where a \textbf{mass} 质量 feels a \textbf{force} 力 from other masses. The \textbf{gravitational field strength} 重力场强度 $g$ at a point is the \textbf{gravitational force per unit mass} on a small \textbf{test mass} 检验质量 placed there:

\[
g = \frac{F}{m}.
\]

Unit: $\text{N kg}^{-1}$ (the same as $\text{m s}^{-2}$ — the acceleration of free fall in the field). $g$ is a \textbf{vector} 矢量, pointing the way the force acts — towards the source mass.

\subsection*{Field lines}

A gravitational field is drawn with \textbf{field lines} 场线 that point the way the force acts on a test mass:

\begin{itemize}
\item around a \textbf{point mass} 质点 or a uniform sphere (treated as a point mass from outside), the field lines are \textbf{radial} 径向, pointing \textbf{inwards}.

\item near the Earth's surface over a small area, the field lines are nearly \textbf{parallel and equally spaced}, pointing straight down — a \textbf{uniform field} 匀强场.

\end{itemize}

Closer lines mean a stronger field.

\par\medskip
\noindent\begin{minipage}{\linewidth}
\centering
\adjustbox{max width=\linewidth, max totalheight=0.42\textheight}{%
  \includegraphics{/Users/joschka/Desktop/Programming/eXercise/output/eXam/handouts/a_level_physics/assets/13-p248-field-lines-strength.pdf}}\\[0.3em]
\emph{Field-line spacing shows the field strength — closer lines mean a stronger field}
\end{minipage}\par\medskip

\section*{Newton's law of gravitation}

For two point masses $m_{1}, m_{2}$ a distance $r$ apart, the force on each is

\[
F = \frac{G m_{1} m_{2}}{r^{2}},
\]

pulling them together along the line joining them. This is \textbf{Newton's law of gravitation} 万有引力定律. The constant $G = 6.67 \times 10^{-11}\ \text{N m}^{2}\ \text{kg}^{-2}$ is the \textbf{universal gravitational constant} 万有引力常量.

\subsection*{Spheres treated as point masses}

For a uniform sphere (such as a planet or star), the field at any point \textbf{outside} is the same as that of a point mass equal to the total mass at the centre. So from above the surface, you can treat the Earth as a point mass at its centre. (Points inside a sphere are different, and are not in the syllabus.)

\par\medskip
\noindent\begin{minipage}{\linewidth}
\centering
\adjustbox{max width=\linewidth, max totalheight=0.42\textheight}{%
  \includegraphics{/Users/joschka/Desktop/Programming/eXercise/output/eXam/handouts/a_level_physics/assets/13-p249-radial-field-sphere.pdf}}\\[0.3em]
\emph{Outside a uniform sphere the field is radial, exactly like a point mass at the centre}
\end{minipage}\par\medskip

\subsection*{Field strength from a point mass}

Put the gravitational force on a test mass $m$ at distance $r$ from a point mass $M$ into $g = F/m$:

\[
F = \frac{G M m}{r^{2}}, \qquad g = \frac{G M}{r^{2}}.
\]

So $g$ falls off as $1/r^{2}$ as you move away from the source.

\subsection*{Why $g$ is nearly constant near the Earth's surface}

The Earth's radius is $R \approx 6.4 \times 10^{6}\ \text{m}$. Rising to height $h$ changes the distance from the centre from $R$ to $R + h$. For $h \ll R$ (any building or mountain), $(R + h)/R \approx 1$, so $g$ barely changes — going from $5\ \text{m}$ to $10\ \text{m}$ high changes $r$ by about one part in a million. In the laboratory, $g$ is effectively constant.

\section*{Orbital motion in a gravitational field}

\par\medskip
\noindent\begin{minipage}{\linewidth}
\centering
\adjustbox{max width=\linewidth, max totalheight=0.42\textheight}{%
  \includegraphics{/Users/joschka/Desktop/Programming/eXercise/output/eXam/handouts/a_level_physics/assets/13-p250-planet-orbit.pdf}}\\[0.3em]
\emph{Gravity provides the centripetal force that keeps a planet in a circular orbit}
\end{minipage}\par\medskip

For a \textbf{satellite} 卫星 of mass $m$ in a circular \textbf{orbit} 轨道 of radius $r$ around a body of mass $M$, gravity provides the \textbf{centripetal force} 向心力:

\[
\frac{G M m}{r^{2}} = \frac{m v^{2}}{r}.
\]

Cancel $m$ (the orbital speed does not depend on the satellite's mass):

\[
v = \sqrt{\frac{G M}{r}}.
\]

The \textbf{period} 周期 follows from $T = 2\pi r / v$:

\[
T = 2\pi \sqrt{\frac{r^{3}}{G M}}, \quad\text{so}\quad T^{2} = \frac{4\pi^{2}}{G M} \cdot r^{3}.
\]

This is \textbf{Kepler's third law} 开普勒第三定律 for circular orbits: $T^{2} \propto r^{3}$. A plot of $T^{2}$ against $r^{3}$ is a straight line through the origin with gradient $4\pi^{2}/(GM)$, so orbital data gives the central mass.

\subsection*{Geostationary orbit}

A \textbf{geostationary} 地球同步 satellite:

\begin{itemize}
\item stays directly above the same point on the Earth (so a fixed dish always points at it),

\item has a period of \textbf{24 hours} (the same as the Earth's rotation),

\item orbits \textbf{west to east} (the same way the Earth turns),

\item must be directly above the \textbf{equator} 赤道.

\end{itemize}

It must have the same \textbf{angular speed} 角速度 as the Earth, in the same direction, in the equatorial plane (or it would drift north–south during the day). From $T = 24\ \text{h}$ and $T^{2} = 4\pi^{2} r^{3}/(GM)$, the radius is $r \approx 4.2 \times 10^{7}\ \text{m}$ (about $3.6 \times 10^{7}\ \text{m}$ above the surface).

\section*{Gravitational potential}

\textbf{Gravitational potential} 引力势 $\phi$ at a point is the \textbf{work done per unit mass} in bringing a small test mass from \textbf{infinity} 无穷远 to that point:

\[
\phi = \frac{W}{m}.
\]

Unit: $\text{J kg}^{-1}$.

The potential is taken as \textbf{zero at infinity}. As the test mass falls in towards the source, gravity does the work for you, so $\phi$ is \textbf{negative} everywhere except at infinity. For a point mass $M$ at distance $r$:

\[
\phi = -\frac{G M}{r}.
\]

$\phi$ is a \textbf{scalar} 标量. For several masses, add the potentials.

\subsection*{Gravitational potential energy of two point masses}

If a test mass $m$ sits where the potential is $\phi$, the \textbf{gravitational potential energy} 重力势能 of the pair is

\[
E_{\text{P}} = m \phi = -\frac{G M m}{r}.
\]

Like the potential, $E_{\text{P}}$ is \textbf{negative} and reaches zero only at infinite separation. Closer masses have more negative potential energy (more tightly bound).

\subsection*{Link with $\Delta E_{\text{P}} = mg\Delta h$}

For small height changes near the surface, $r$ barely changes, so $\Delta E_{\text{P}} \approx mg\Delta h$. For large changes (a satellite moving to a higher orbit) use $-GMm/r$ at each radius and take the difference:

\[
\Delta E_{\text{P}} = GMm\left(\frac{1}{r_{1}} - \frac{1}{r_{2}}\right) \quad (r_{2} > r_{1}),
\]

which is positive (energy must be supplied to raise the satellite).

\subsection*{Escape velocity (from conservation of energy)}

To escape from radius $r$ to infinity, an object's \textbf{kinetic energy} 动能 must equal the size of its gravitational potential energy:

\[
\tfrac{1}{2} m v_{\text{esc}}^{2} = \frac{G M m}{r}, \qquad v_{\text{esc}} = \sqrt{\frac{2 G M}{r}}.
\]

At the Earth's surface, the \textbf{escape velocity} 逃逸速度 is $\approx 11\ \text{km s}^{-1}$. It does not depend on the object's mass.


\end{document}
